\clearpage\pagestyle{otMain}\fontsize{11}{15}\selectfont\corrigesChapitre{11}{Probabilités et expériences à deux étapes}
\begin{blocLivre}{cor-11-1}\begin{corrige}{1}{Un dé}\demarche{Énumérer les issues favorables.}
Les issues favorables sont 2, 4 et 6, parmi six issues équiprobables. La probabilité vaut $3/6=\res{1/2}$.
\end{corrige}\end{blocLivre}
\begin{blocLivre}{cor-11-2}\begin{corrige}{2}{Des boules colorées}\demarche{Compter les boules plutôt que les couleurs.}
Le total vaut 10. $P(R)=5/10=\res{1/2}$ et $P(\text{non }R)=1-1/2=\res{1/2}$. Trois couleurs ne signifient pas trois issues équiprobables.
\end{corrige}\end{blocLivre}
\begin{blocLivre}{cor-11-3}\begin{corrige}{3}{Impossible ou certain}\demarche{Comparer l’événement à toutes les issues.}
Obtenir 9 est impossible : probabilité $\res0$. Toutes les faces sont inférieures à 7 : probabilité $\res1$.
\end{corrige}\end{blocLivre}
\begin{blocLivre}{cor-11-4}\begin{corrige}{4}{Une fréquence expérimentale}\demarche{Ne pas confondre résultat observé et modèle.}
La fréquence est $46/100=\res{0{,}46}$. Cet écart à 0,5 peut se produire avec une pièce équilibrée ; ce seul résultat n’est pas une preuve. La fréquence dépend de l’expérience réalisée.
\end{corrige}\end{blocLivre}
\begin{blocLivre}{cor-11-5}\begin{corrige}{5}{Une roue à secteurs égaux}\demarche{Identifier les secteurs équiprobables.}
Les numéros favorables sont 3 et 6 : $2/8=\res{1/4}$. L’égalité des angles et l’hypothèse du modèle justifient le comptage.
\end{corrige}\end{blocLivre}
\begin{blocLivre}{cor-11-6}\begin{corrige}{6}{Deux pièces}\demarche{Suivre toutes les feuilles de l’arbre.}
Les issues sont PP, PF, FP, FF. Chacune vaut $1/2\times1/2=1/4$. PF et FP conviennent, d’où $\res{1/2}$. L’ordre distingue les deux issues.
\label{fig-cor-11-6}\figureGeo[corrige]{proba-pieces}
\end{corrige}\end{blocLivre}
\begin{blocLivre}{cor-11-7}\begin{corrige}{7}{La somme de deux dés}\demarche{Compter des couples et non des sommes.}
Les 36 couples ordonnés sont équiprobables. Les six couples $(1;6),(2;5),(3;4),(4;3),(5;2),(6;1)$ conviennent : $\res{1/6}$. Les sommes elles-mêmes ne sont pas équiprobables.
\label{fig-cor-11-7}\figureGeo[corrige]{proba-des}
\end{corrige}\end{blocLivre}
\begin{blocLivre}{cor-11-8}\begin{corrige}{8}{Avec remise}\demarche{Tenir compte de la remise avant la deuxième branche.}
La composition reste identique : le chemin RR a pour poids $3/5\times3/5=\res{9/25}$. La remise maintient cinq boules à la seconde étape.
\label{fig-cor-11-8}\figureGeo[corrige]{proba-remise}
\end{corrige}\end{blocLivre}
\begin{blocLivre}{cor-11-9}\begin{corrige}{9}{Sans remise}\demarche{Actualiser la composition après la première boule.}
Après une rouge, il reste 2 rouges sur 4 boules. La probabilité de RR vaut $3/5\times2/4=\res{3/10}$. Elle est différente du cas avec remise.
\label{fig-cor-11-9}\figureGeo[corrige]{proba-sans-remise}
\end{corrige}\end{blocLivre}
\begin{blocLivre}{cor-11-10}\begin{corrige}{10}{Au moins un six}\demarche{Passer au contraire pour éviter un double comptage.}
L’événement contraire est « aucun six », de probabilité $5/6\times5/6=25/36$. La probabilité cherchée vaut $1-25/36=\res{11/36}$. Additionner $1/6+1/6$ aurait compté deux fois le couple $(6;6)$.
\end{corrige}\end{blocLivre}
\begin{blocLivre}{cor-11-11}\begin{corrige}{11}{Deux couleurs différentes}\demarche{Additionner les deux chemins favorables.}
Les chemins RB et BR sont incompatibles. Leurs poids valent $3/5\times2/4=3/10$ et $2/5\times3/4=3/10$. La somme est $\res{3/5}$. Les deux ordres comptent.
\label{fig-cor-11-11}\figureGeo[corrige]{proba-sans-remise}
\end{corrige}\end{blocLivre}
\begin{blocLivre}{cor-11-12}\begin{corrige}{12}{Un jeu à points}\demarche{Distinguer gain brut, mise et bilan net.}
Le gain brut moyen vaut $8/6+4/6=\res2$ points. Il égale la mise : le bilan net moyen théorique est 0. Cela ne garantit ni un remboursement ni un gain à chaque partie ; un résultat individuel peut être $-2$, 2 ou 6 points nets.
\end{corrige}\end{blocLivre}
\begin{blocLivre}{cor-11-13}\begin{corrige}{13}{Un effectif inconnu}\demarche{Traduire le total avant le produit en croix.}
On écrit $6/(6+b)=2/5$. Les produits en croix donnent $30=12+2b$, donc $\res{b=9}$. Avec 15 boules, $6/15=2/5$, ce qui confirme le résultat.
\end{corrige}\end{blocLivre}
\begin{blocLivre}{cor-11-14}\begin{corrige}{14}{Un code à deux chiffres}\demarche{Préciser si répétition et zéro initial sont autorisés.}
Il y a $4\times4=16$ codes équiprobables. Pour chaque premier chiffre, trois seconds sont différents : $4\times3=12$. La probabilité vaut $12/16=\res{3/4}$. Le complément correspond à 00, 11, 22, 33.
\end{corrige}\end{blocLivre}
\begin{blocLivre}{cor-11-15}\begin{corrige}{15}{Estimer sans promettre}\demarche{Comparer les écarts absolus sans en faire une règle monotone.}
Les fréquences sont 0,327 et 0,3354. Leurs écarts à $1/3$ valent environ 0,00633 et 0,00207. La deuxième est plus proche dans ces données. Aucune amélioration à chaque étape n’est garantie : un prochain succès ou échec peut augmenter l’écart.
\end{corrige}\end{blocLivre}
